% Generated draft: requires mathematical and source review.
\documentclass[11pt]{article}
\usepackage[margin=25mm]{geometry}
\usepackage{fontspec}
\setmainfont{Latin Modern Roman}
\usepackage{amsmath,amssymb,mathtools,unicode-math}
\setmathfont{Latin Modern Math}
\usepackage{enumitem,needspace,xurl,hyperref}
\setlength{\parindent}{0pt}
\setlength{\parskip}{5pt}
\setlength{\emergencystretch}{3em}
\newcommand{\N}{\mathbb{N}}
\newcommand{\Z}{\mathbb{Z}}
\newcommand{\Q}{\mathbb{Q}}
\newcommand{\R}{\mathbb{R}}
\newcommand{\C}{\mathbb{C}}
\newcommand{\F}{\mathbb{F}}
\newcommand{\A}{\mathbb{A}}
\newcommand{\PP}{\mathbb{P}}
\newcommand{\E}{\mathbb{E}}
\newcommand{\eps}{\varepsilon}
\newcommand{\dd}{\,\mathrm{d}}
\begin{document}
\section*{UnsolvedMath ID: 30002330}
\textbf{Source problem code:} catalogue entry
\subsection{Definitions and mathematical statement}
Let \(\Sigma=f(S^2)\subset\R^3\) be the image of a smooth embedding of the two-dimensional sphere. Thus \(\Sigma\) is compact and has no boundary; smooth is understood in the customary \(C^\infty\) sense. For a choice of unit normal, let \(A_q:T_q\Sigma\to T_q\Sigma\) denote the shape operator at \(q\). A point \(p\in\Sigma\) is an umbilic if \(A_p=\lambda\operatorname{Id}_{T_p\Sigma}\) for some \(\lambda\in\R\). It is isolated if some disk neighborhood \(U\) of \(p\) contains no other umbilics.

On \(U\setminus\{p\}\), either eigenspace of the shape operator defines a principal curvature line field. To define its index, choose an oriented orthonormal tangent frame over \(U\) and a positively oriented simple loop \(\gamma:[0,1]\to U\setminus\{p\}\) enclosing \(p\) once. Write the chosen, unoriented principal line along \(\gamma\) as an angle modulo \(\pi\), and let \(\theta:[0,1]\to\R\) be a continuous lift of that angle. Set
\[
\operatorname{ind}_p(\Sigma)=\frac{\theta(1)-\theta(0)}{2\pi}\in\tfrac12\Z.
\]
This index is independent of the frame extending across \(U\), of a sufficiently small enclosing loop, and of which of the two orthogonal principal line fields is chosen.

The supplied question is whether
\[
\exists\,f:S^2\hookrightarrow\R^3\text{ smooth},\quad
\exists\,p\in\Sigma=f(S^2)
\]
such that \(\Sigma\) is strictly convex, \(p\) is an isolated umbilic, and
\[
\operatorname{ind}_p(\Sigma)=\frac32.
\]
Here the index concerns an unoriented principal line field, so the required change of its lifted angle around \(p\) is \(3\pi\).

The supplied record does not explicitly define strict convexity. Two conventions that require distinction are: \(\Sigma=\partial K\) for a compact convex body \(K\subset\R^3\) with nonempty interior such that every open segment joining distinct points of \(K\) lies in \(\operatorname{int}K\); and definiteness of the second fundamental form at every point of \(\Sigma\), with a consistent normal and sign convention. The former permits vanishing principal curvature and does not by itself impose the latter. The intended convention must be checked against the source; this draft retains the source wording rather than choosing between these hypotheses.
\subsection{Short English statement}
Does a smooth, strictly convex embedded sphere in \(\R^3\) have an isolated umbilic at which the principal curvature line field has index \(3/2\)? The precise meaning of strict convexity requires source review. The supplied registry labels the problem open; that status is not independently verified here.
\subsection{Sources}
\begin{itemize}
\item Supplied local UnsolvedMath registry record, permanent ID \(30002330\), exact source code \texttt{catalogue entry}, title: \emph{Convex Spheres with an Isolated Index-Three-Halves Umbilic}. Collection: Oberwolfach Reports Open Problems. The supplied proposed year is 2013; no proposer is identified.
\item Supplied source citation: \emph{Progress in Surface Theory (2014)}, Oberwolfach Reports, pp.\ 1253--1312, DOI \texttt{10.4171/owr/2013/21}. \url{https://doi.org/10.4171/owr/2013/21}. The supplied background describes these reports as records of mathematical workshops and their open-problem discussions. The citation's 2014 label and the supplied 2013 year are retained without reconciliation.
\item Guilfoyle--Klingenberg, \emph{An index bound for smooth umbilic points}. Supplied literature link: \url{https://arxiv.org/abs/1207.5994}.
\item \emph{Roots of Polynomials and Umbilics of Surfaces}. Supplied literature link: \url{https://doi.org/10.1007/s00025-023-02003-4}.
\item Supplied literature assessment, checked at \texttt{2026-08-21}: the bound \(\operatorname{Index}\leq 3/2\) is described as known, existence of a smooth convex umbilic of index \(3/2\) as unproved, and no counterexample as found in that review. This is supplied dated curation derived from the OWR record, not an independently verified literature assessment. The hypotheses and scope of the cited index bound have not been established here.
\item The supplied statement status is \texttt{corrected\_verified}; its accompanying verification note reports comparison with the official Oberwolfach report and normalization of notation, context, and question form without changing the source-level claim. This is supplied provenance, not verification performed for this draft.
\end{itemize}
\end{document}
